\section{Scalar regularity and properties of a minimizer}\label{static-sec:scalar}

We now establish the scalar identities used in the overlap comparisons
and the equilibrium energy estimate. The argument has three parts.
First, we justify the scalar equation and its stochastic evolution for
general distribution functions, including those with atoms. Next, we
derive first variations and use them to locate the support of an
arbitrary minimizing measure. Finally, the resulting stationarity
identity gives a strict energy bound when $\beta>1$. No further
description of the minimizer's support is assumed.

Regularity and step approximation for the Parisi evolution, together
with zero-field support properties of its minimizers, were developed
by Auffinger and Chen~\cite[Propositions~1--2 and Theorems~1, 5]{AC2015}.
The full first-variation characterization
appears in Chen~\cite[Theorem~2 and Proposition~1]{Chen2017} and,
independently, in Jagannath and Tobasco~\cite[Proposition~1.1,
Corollary~3.10 and Remark~1.3]{JT2017}. We prove the required statements
here with the normalization and atom conventions used by the later
pressure comparisons.

Throughout this section, $\alpha(s)=\mu([0,s])$ is the right-continuous
distribution function of a probability measure on $[0,1]$. Recall the
scalar terminal-value equation
\begin{equation}\label{static-eq:scalar-pde-recalled}
 -\partial_s\Phi_\alpha
 =\frac{\beta^2}{2}\bigl(\partial_{zz}\Phi_\alpha
          +\alpha(s)(\partial_z\Phi_\alpha)^2\bigr),
 \qquad \Phi_\alpha(1,z)=\log(2\cosh z).
\end{equation}
For a step distribution function this is the Gaussian recursion from
the preceding section. For general distribution functions, the scalar
values are defined by approximation, which the next lemma justifies.
The time derivative in the equation is understood almost everywhere;
no differentiability at a jump of $\alpha$ is required.
We denote the first two spatial derivatives by $u=\partial_z\Phi_\alpha$ and
$\chi=\partial_{zz}\Phi_\alpha$.

\subsection{Regularity and first variations}

\begin{lemma}[Regularity, stability, and scalar evolution]
\label{static-lem:scalar-regularity}
The step approximation defining $\Phi_\alpha$ is independent of the
approximating sequence. The functions $\Phi_\alpha,u,\chi$ are continuous
on $[0,1]\times\R$, with
\begin{equation}\label{static-eq:scalar-derivative-bounds}
 |u|\le1,\qquad 0<\chi\le1,\qquad
 |\chi(s,z)-\chi(s,z')|\le C_\beta|z-z'|.
\end{equation}
The function $\Phi_\alpha(s,\cdot)$ is even, so $u(s,0)=0$.
For two distribution functions $\alpha,\widetilde\alpha$,
\begin{equation}\label{static-eq:scalar-stability}
 \sup_{s,z}\bigl(
 |\Phi_\alpha-\Phi_{\widetilde\alpha}|
 +|u_\alpha-u_{\widetilde\alpha}|
 +|\chi_\alpha-\chi_{\widetilde\alpha}|
 \bigr)
 \le C_\beta\int_0^1|\alpha(s)-\widetilde\alpha(s)|\dd s.
\end{equation}
These functions are also continuous in $\beta>0$, uniformly in $(s,z)$
for their differences, locally in $\beta$.

The stochastic differential equation
\begin{equation}\label{static-eq:scalar-diffusion}
 \dd X_s=\beta\dd B_s+\beta^2\alpha(s)u(s,X_s)\dd s,
 \qquad X_0=0,
\end{equation}
has a unique strong solution. For a step distribution function, its
transitions across the constant-mass intervals are the normalized tilted
Gaussian kernels of the scalar recursion. Its bounded martingale
$U_s=u(s,X_s)$ satisfies
\begin{equation}\label{static-eq:scalar-martingale}
 U_t-U_s=\beta\int_s^t\chi(v,X_v)\dd B_v.
\end{equation}
Moreover,
\begin{equation}\label{static-eq:scalar-chi-expectation}
 \E\chi(t,X_t)-\E\chi(s,X_s)
 =-\beta^2\int_s^t\alpha(v)\E\chi(v,X_v)^2\dd v.
\end{equation}
For $0\le s<t\le1$, the conditional form is
\begin{equation}\label{eq:scalar-evolution}
 \E[\chi(t,X_t)\mid X_s]
 =\chi(s,X_s)-\beta^2\int_s^t\alpha(v)
                   \E[\chi(v,X_v)^2\mid X_s]\dd v.
\end{equation}
Consequently $\Gamma(s):=\E u(s,X_s)^2$ is continuously differentiable,
with one-sided derivatives at the endpoints, and
\begin{equation}\label{static-eq:gamma-derivative}
 \Gamma(0)=0,\qquad
 \Gamma'(s)=\beta^2\E\chi(s,X_s)^2.
\end{equation}
The well-posedness statement and the evolution identities
\eqref{static-eq:scalar-martingale}--\eqref{static-eq:scalar-chi-expectation} also hold
for a diffusion started at any deterministic time and spatial value.
\end{lemma}

\begin{proof}
We prove estimates uniform over step distribution functions, use them
to pass to general $\alpha$, and then establish continuity in the
temperature parameter.

\par\medskip\noindent\textbf{Derivative bounds along a finite step path.}\enspace
Suppose first that $\alpha$ is a step function. On an interval of
constant mass $m$, its recursion is
\[
 T_m f(z)=m^{-1}\log\E e^{m f(z+\beta\sqrt{\ell}\,Z)}
\]
where $\ell$ is the interval length and
$T_0f=\E f(z+\beta\sqrt{\ell}\,Z)$. To keep track of the successive
tilts, fix $(s,z)$ and list the remaining breakpoints as
$s=t_0<\cdots<t_r=1$. Let $m_j$ be the mass on
$[t_j,t_{j+1})$ and set $f_j=\Phi_\alpha(t_j,\cdot)$.
Starting with $Z_0=z$, sample $Z_{j+1}$ from
\[
 K_j(v,\dd y)
 =\exp\{m_j[f_{j+1}(y)-f_j(v)]\}
       N(v,\beta^2(t_{j+1}-t_j))(\dd y),
\]
where the exponential equals one at $m_j=0$. Each kernel is normalized by
the scalar recursion on its own interval. Put
$\mathcal F_j=\sigma(Z_0,\ldots,Z_j)$ and $M_j=f_j'(Z_j)$.
Differentiating the recursion once and twice gives
\begin{align*}
 \E[M_{j+1}\mid\mathcal F_j]&=M_j,\\
 f_j''(Z_j)&=\E[f_{j+1}''(Z_{j+1})\mid\mathcal F_j]
      +m_j\E[(M_{j+1}-M_j)^2\mid\mathcal F_j].
\end{align*}
Thus $M$ is a martingale for the concatenated kernels, even though
the tilts on successive intervals can differ. Its terminal value is
$M_r=\tanh Z_r$. Iterating the second identity yields
\[
 \chi(s,z)=\E\operatorname{sech}^2(Z_r)
           +\sum_{j=0}^{r-1}m_j\E(M_{j+1}-M_j)^2.
\]
The increments are orthogonal, so their unweighted squared sum has
expectation $\Var(\tanh Z_r)$. Since $0\le m_j\le1$, we obtain
\[
 0<\E\bigl[1-\tanh^2(Z_r)\bigr]
 \le\chi(s,z)
 \le \E\bigl[1-\tanh^2(Z_r)\bigr]
       +\Var(\tanh Z_r)
 =1-u(s,z)^2\le1.
\]
Every recursion preserves evenness. Its gradient is a conditional
mean of the terminal gradient, which gives $|u|\le1$.

Before differentiating the equation, we record why all required
spatial derivatives exist and are bounded for a fixed finite path.
The $k$th derivative of one logarithmic Gaussian integral is a
polynomial in tilted moments of $f',\ldots,f^{(k)}$, with coefficients
bounded for $0\le m\le1$. At $m=0$, differentiation passes directly
through Gaussian expectation. All positive-order derivatives of
$\log(2\cosh z)$ are bounded, so finite induction gives bounded
derivatives at the breakpoints and within each interval. The scalar
values have at most linear growth, ensuring Gaussian integrability
at every step. These preliminary bounds may depend on the path; we
next obtain the uniform bound needed for approximation.

Within a constant-mass interval the recursion is smooth and solves
Equation~\eqref{static-eq:scalar-pde-recalled}. Its first two spatial
derivatives satisfy
\begin{align*}
 -u_s&=\frac{\beta^2}{2}u_{zz}+\beta^2\alpha u u_z,\\
 -\chi_s&=\frac{\beta^2}{2}\chi_{zz}
       +\beta^2\alpha u\chi_z+\beta^2\alpha\chi^2.
\end{align*}
Set $w=\chi_z$. One more differentiation gives a linear equation:
\[
 -w_s=\frac{\beta^2}{2}w_{zz}
       +\beta^2\alpha u w_z+3\beta^2\alpha\chi w.
\]
The bounds already proved for $u$ and $\chi$ control this equation's
drift and potential uniformly over all step paths. Its terminal value
has absolute value at most $2$. Concatenate the diffusion
representations over the finitely many intervals to obtain
\[
 \sup_{s,z}|w(s,z)|\le2e^{3\beta^2}.
\]
These differentiations may all be performed first inside the finite
Gaussian recursions. Their bounded derivatives justify passage under
the integrals, so the uniform estimate does not presuppose any
regularity of the limiting path.

\par\medskip\noindent\textbf{The scalar diffusion and its evolution identities.}\enspace
We next identify the tilted kernels with the diffusion in
Equation~\eqref{static-eq:scalar-diffusion}. On a constant-mass
interval $[a,b]$, start Brownian motion at $z$ by setting
$Y_s=z+\beta(B_s-B_a)$. Under this law, It\^o's formula identifies
\[
 \exp\{m(\Phi(b,Y_b)-\Phi(a,z))\}
\]
as the exponential martingale with stochastic logarithm
$m\beta\int_a^b u(v,Y_v)\dd B_v$. Its bounded integrand justifies
the change of measure, which adds drift $\beta^2m u$. Under that
drift, a further application of It\^o's formula gives
\[
 \dd u(s,X_s)=\beta\chi(s,X_s)\dd B_s,
 \qquad
 \dd\chi(s,X_s)=\beta\chi_z(s,X_s)\dd B_s
                 -\beta^2\alpha(s)\chi(s,X_s)^2\dd s.
\]
This establishes both evolution identities for step paths.

\par\medskip\noindent\textbf{Stability and passage to general distribution functions.}\enspace
The remaining task is to make these constructions independent of the
step approximation. Subtract the equations for two step paths, keeping
the same diffusion coefficient. The difference of the scalar values
has drift $\frac{\beta^2}{2}\alpha(u+\widetilde u)$ and source
$\frac{\beta^2}{2}(\alpha-\widetilde\alpha)\widetilde u^2$.
The equation for $u-\widetilde u$ has bounded drift and potential,
with source bounded by $C_\beta|\alpha-\widetilde\alpha|$.
For $\chi-\widetilde\chi$, the drift and potential remain bounded,
and the source is bounded by
\[
 C_\beta\bigl(|\alpha-\widetilde\alpha|(s)
                 +\sup_z|u(s,z)-\widetilde u(s,z)|\bigr).
\]
To obtain the last bound, expand
$\alpha u\chi_z-\widetilde\alpha\widetilde u\widetilde\chi_z$
and use the uniform bound on $\widetilde\chi_z$. In the remaining
reaction term, factor the difference of squares. Each difference
has zero terminal value, so the diffusion representation of these
linear equations proves Equation~\eqref{static-eq:scalar-stability},
first for $\Phi$, then for $u$, and finally for $\chi$.

Approximate an arbitrary distribution function in $L^1$ by step
distribution functions. The stability estimates make $\Phi,u,\chi$
converge uniformly, with uniform convergence of differences for the
unbounded scalar value $\Phi$. The fundamental theorem of calculus
in $z$ identifies the limits of $u$ and $\chi$ with the first two
spatial derivatives of the limiting $\Phi$. Continuity, symmetry,
and all non-strict bounds follow by passage to the limit.
Passing to the limit in the time-integrated step equation also gives
Equation~\eqref{static-eq:scalar-pde-recalled} almost everywhere in time.

The limiting drift in Equation~\eqref{static-eq:scalar-diffusion} is
measurable in time, bounded, and uniformly Lipschitz in space. Hence
the equation has a strong solution and pathwise uniqueness. Couple
the solutions for two step approximations with the same Brownian
motion; Gronwall's inequality gives convergence uniformly in time.
We may therefore also pass the terminal lower bound for $\chi$ to
the limit:
\[
 \chi(s,z)\ge\E_{s,z}\operatorname{sech}^2(X_1)>0.
\]
Since $X_1$ is finite almost surely, this lower bound is strictly
positive. It\^o isometry passes the martingale identity to the limit,
and bounded convergence does the same for the expectation identity.
Nothing in these arguments requires the starting point to be
$(0,0)$; they apply at every deterministic starting time and spatial
value. Applying the expectation identity to a diffusion started at
$(s,z)$ and then using the Markov property gives
Equation~\eqref{eq:scalar-evolution}. Finally,
Equation~\eqref{static-eq:scalar-martingale} and the
continuity of $\chi$ and $X$ give
Equation~\eqref{static-eq:gamma-derivative}.

\par\medskip\noindent\textbf{Continuity in the temperature parameter.}\enspace
Write $b=\beta^2$ and repeat the subtraction argument with different
values of $b$. In the equation for $\Phi$, the additional source is
$(\widetilde b-b)(\widetilde\chi+
\alpha\widetilde u^2)/2$, bounded by $C|\widetilde b-b|$.
The bound on $\widetilde\chi_z$ gives the same source bound for
$u$. Thus the differences of $\Phi$ and $u$ are
$O(|\widetilde b-b|)$ locally uniformly for $b>0$.
To obtain continuity of $\chi$, use its uniform spatial Lipschitz
bound and approximate it by a difference quotient of $u$. For
any $h>0$,
\[
 \sup_{s,z}|\chi-\widetilde\chi|
 \le\frac{2}{h}\sup_{s,z}|u-\widetilde u|+C h.
\]
Choose $h$ small first, then take $\widetilde b$ sufficiently close to
$b$. This proves the required uniform continuity of $\chi$.
The step approximation transfers these parameter estimates to
all distribution functions.
\end{proof}

We next differentiate the scalar functional in two directions: a
change of its distribution function and a change of temperature with
the distribution function fixed. We prove both formulas by differences,
so they require no differentiability of a minimizing measure.
For another distribution function $\widetilde\alpha$, set
$\alpha_\theta=\alpha+\theta(\widetilde\alpha-\alpha)$. The first
variation is
\begin{equation}\label{static-eq:scalar-mass-variation}
 \left.\frac{\dd}{\dd\theta}\cP_\beta(\alpha_\theta)
       \right|_{\theta=0+}
 =\frac{\beta^2}{2}\int_0^1
       (\widetilde\alpha-\alpha)(s)(\Gamma(s)-s)\dd s.
\end{equation}
We use $\cP_\beta(\alpha)$ to denote the functional of the measure
having distribution function $\alpha$. With $\alpha$ held fixed,
the temperature derivative of the scalar value is
\begin{equation}\label{static-eq:scalar-temperature-variation}
 \frac{\partial}{\partial(\beta^2)}\Phi_\alpha(0,0)
 =\frac12\int_0^1
     \E\bigl[\chi(s,X_s)+\alpha(s)u(s,X_s)^2\bigr]\dd s.
\end{equation}
For each formula, start with step paths, where the equations are
classical between finitely many breakpoints. Subtraction gives an
integral representation, and that representation passes to general
paths by the stability lemma.

\par\medskip\noindent\textbf{Variation of the distribution function.}\enspace
Subtract the equations for $\Phi_{\alpha_\theta}$ and $\Phi_\alpha$.
The resulting drift is
$\beta^2\alpha(u_{\alpha_\theta}+u_\alpha)/2$, and the source is
$\theta\beta^2(\widetilde\alpha-\alpha)u_{\alpha_\theta}^2/2$.
Introduce $Y^\theta$ by $Y^\theta_0=0$ and
\[
 \dd Y_s^\theta=\beta\dd B_s
 +\frac{\beta^2}{2}\alpha(s)
       (u_{\alpha_\theta}+u_\alpha)(s,Y_s^\theta)\dd s.
\]
The diffusion representation of the difference equation is the
exact identity
\begin{equation}\label{static-eq:scalar-mass-difference}
 \frac{\Phi_{\alpha_\theta}(0,0)-\Phi_\alpha(0,0)}{\theta}
 =\frac{\beta^2}{2}\E\int_0^1
       (\widetilde\alpha-\alpha)(s)
       u_{\alpha_\theta}(s,Y_s^\theta)^2\dd s.
\end{equation}
For general $\alpha,\widetilde\alpha$, choose step approximations
converging to each in $L^1$, and take their convex combinations at
a fixed $\theta>0$. Equation~\eqref{static-eq:scalar-stability}
shows that the mixed drifts converge in integrated spatial supremum
norm. They also have a common spatial Lipschitz bound. A coupling
with the same Brownian motion and Gronwall's inequality therefore
gives uniform convergence of the corresponding $Y^\theta$
processes. The sources converge in integrated expectation by
$|u|\le1$, the spatial Lipschitz bound for $u$, its uniform
convergence, and the $L^1$ convergence of
$\widetilde\alpha-\alpha$. Hence
Equation~\eqref{static-eq:scalar-mass-difference} holds for arbitrary
distribution functions.

We may now send $\theta\downarrow0$. The same estimates give
$Y^\theta\to X$ and uniform convergence of $u_{\alpha_\theta}$ to
$u_\alpha$. The right-hand side converges to
$\frac{\beta^2}{2}\int(\widetilde\alpha-\alpha)\Gamma$.
Subtract the derivative of the linear penalty to obtain
Equation~\eqref{static-eq:scalar-mass-variation}.

\par\medskip\noindent\textbf{Variation of the temperature.}\enspace
Set $b=\beta^2$ and use a superscript $b$ for scalar functions with
parameter $\sqrt b$. For $\widetilde b\ne b$, let
$Y^{b,\widetilde b}_0=0$ and define its evolution by
\[
 \dd Y_s^{b,\widetilde b}=\sqrt b\dd B_s
 +\frac b2\alpha(s)(u^b+u^{\widetilde b})
                (s,Y_s^{b,\widetilde b})\dd s.
\]
Subtract the two step-path equations with diffusion coefficient
$b/2$ to obtain
\begin{equation}\label{static-eq:scalar-temperature-difference}
 \frac{\Phi^{\widetilde b}(0,0)-\Phi^b(0,0)}{\widetilde b-b}
 =\frac12\E\int_0^1
       [\chi^{\widetilde b}+\alpha(u^{\widetilde b})^2]
                    (s,Y_s^{b,\widetilde b})\dd s.
\end{equation}
For fixed $b,\widetilde b$, pass to general $\alpha$ by the same
step approximation and Brownian coupling used above. The additional
term involving $\chi$ is controlled by its uniform spatial
Lipschitz bound. Next let $\widetilde b\to b$. Parameter continuity
in Lemma~\ref{static-lem:scalar-regularity} and the coupling give
$Y^{b,\widetilde b}\to X$. The integrands remain bounded and
converge in integrated expectation. The limit in
Equation~\eqref{static-eq:scalar-temperature-difference} therefore
proves Equation~\eqref{static-eq:scalar-temperature-variation} for
approach from either side.

\subsection{Minimizers and the energy gap}

The mass variation converts minimization into a condition on the
support. We will use that condition both to identify the expected
scalar susceptibility at every time and to show that the minimizing
measure is nontrivial when $\beta>1$.

\begin{proposition}[Support and stationarity identities]
\label{static-prop:scalar-minimizer}
The scalar functional has a minimizer $\mu$. For every such minimizer,
with its associated scalar process, the following assertions hold:
\begin{enumerate}[label=\textup{(\roman*)}]
 \item $0\in\supp\mu$, and $\supp\mu\subset[0,1-\eps_0]$ for some
       $\eps_0>0$.
 \item At every $q\in\supp\mu\cap(0,1)$,
       $\Gamma(q)=q$ and $\Gamma'(q)\le1$. At $q=0$,
       $\Gamma(0)=0$ and $\Gamma'(0)\le1$.
 \item Define
 \begin{equation}\label{static-eq:scalar-A-d-definitions}
   d_*(s)=1-\int_{(s,1]}t\,\mu(\dd t),\qquad
   A(s)=\int_s^1\alpha(t)\dd t=d_*(s)-s\alpha(s).
 \end{equation}
 Then, at every $s\in[0,1]$,
 \begin{equation}\label{static-eq:scalar-stationarity-identity}
   \E\bigl[\chi(s,X_s)+\alpha(s)u(s,X_s)^2\bigr]=d_*(s).
 \end{equation}
 Equivalently, since $d_*(s)=A(s)+s\alpha(s)$,
 \begin{equation}\label{eq:scalar-stationarity}
   \E\chi(s,X_s)=A(s)+\alpha(s)(s-\Gamma(s))
       \qquad(0\le s\le1).
 \end{equation}
 In particular
 \begin{equation}\label{static-eq:scalar-susceptibility}
   \chi(0,0)=A(0),\qquad \beta A(0)\le1.
 \end{equation}
 Furthermore $\alpha(s)>0$ for every $s\in(0,1]$,
 $A(h)<A(0)$ for every $h\in(0,1]$, and $d_*(0+)=A(0)$.
\end{enumerate}
\end{proposition}

\begin{proof}
\par\medskip\noindent\textbf{Existence and variational support condition.}\enspace
Weak convergence of probability measures on $[0,1]$ implies
pointwise convergence of their distribution functions at continuity
points of the limit. Bounded convergence then gives convergence in
$L^1$. By Lemma~\ref{static-lem:scalar-regularity}, the functional
$\cP_\beta$ is continuous for this topology. The space of probability
measures on $[0,1]$ is compact, so a minimizer exists.

Fix any minimizer and define
\[
 G(q)=\int_q^1(\Gamma(s)-s)\dd s.
\]
Fubini's theorem and Equation~\eqref{static-eq:scalar-mass-variation}
give $\int G\dd\widetilde\mu\ge\int G\dd\mu$ for every
probability measure $\widetilde\mu$. Testing with point masses shows
that $\mu$ is supported on the set of global minima of $G$.
The function $G$ is twice continuously differentiable and satisfies
\[
 G'(q)=q-\Gamma(q),\qquad G''(q)=1-\Gamma'(q).
\]
The first- and second-derivative tests now give the claimed identities
and inequalities at interior support points. At the right endpoint,
$\Gamma(1)=\E\tanh^2(X_1)<1$, and hence $G'(1)>0$.
Thus $1$ does not minimize $G$. Its minimum set is compact, so
$\supp\mu$ stays a positive distance below $1$.

\par\medskip\noindent\textbf{The left endpoint of the support.}\enspace
Let $q_0=\min\supp\mu$. We claim that $q_0=0$. Suppose instead
that $q_0\in(0,1)$. Then $\alpha=0$ on $[0,q_0)$, so
$X_s=\beta B_s$ there and the equation for $\chi$ is the backward
heat equation. Consequently $\chi(s,X_s)$ is a bounded martingale
on $0\le s\le q_0$. One can see this either from the heat
semigroup or by choosing step approximations identically zero below
$q_0$. Its second moment is nondecreasing, as is $\Gamma'(s)$.
But $\Gamma(q_0)=q_0$, $\Gamma(0)=0$, and
$\Gamma'(q_0)\le1$. Thus the average of $\Gamma'$ on this
interval is $1$, whereas its maximum is at most $1$. Continuity
forces $\Gamma'\equiv1$ on $[0,q_0]$.

To contradict this constancy, observe that $\chi(q_0,\cdot)$ is
positive everywhere and cannot be constant: otherwise
$u(q_0,\cdot)$ would have a positive constant derivative, contrary
to $|u|\le1$. The function $\chi(q_0,\cdot)$ is continuous, and
$X_{q_0}$ is a nondegenerate Gaussian with full support. Therefore
$\Var(\chi(q_0,X_{q_0}))>0$, while $\chi(0,X_0)$ is
deterministic. The martingale's second moment strictly increases
between these endpoints, which contradicts the constancy of
$\Gamma'$. Hence $q_0=0$. At this endpoint, $G$ has a minimum
and symmetry gives $G'(0)=0$. The one-sided second-derivative test
then yields $\Gamma'(0)\le1$.

\par\medskip\noindent\textbf{Stationarity with atoms included.}\enspace
We have located the support and established $\Gamma(t)=t$ on it.
It remains to turn these support identities into an identity at every
time. Use Equation~\eqref{static-eq:scalar-chi-expectation}, the
terminal relation $\chi(1,z)=1-u(1,z)^2$, and
Equation~\eqref{static-eq:gamma-derivative} to write
\[
 \E\chi(s,X_s)=1-\Gamma(1)+\int_s^1\alpha(t)\Gamma'(t)\dd t.
\]
Stieltjes integration by parts with the right-continuous
$\alpha$ gives the following formula, with every atom retained:
\[
 \int_s^1\alpha(t)\Gamma'(t)\dd t
 =\Gamma(1)-\alpha(s)\Gamma(s)
     -\int_{(s,1]}\Gamma(t)\,\mu(\dd t).
\]
In particular, the integral uses $(s,1]$, while the endpoint term
uses $\alpha(s)$. Substituting $\Gamma(t)=t$ on $\supp\mu$,
including at zero, proves
Equation~\eqref{static-eq:scalar-stationarity-identity}. Fubini's
theorem gives the expression for $A$ in
Equation~\eqref{static-eq:scalar-A-d-definitions}. At $s=0$ the
stationarity identity becomes $\chi(0,0)=A(0)$. Finally,
$\Gamma'(0)=\beta^2\chi(0,0)^2\le1$ proves
Equation~\eqref{static-eq:scalar-susceptibility}.

Because $0\in\supp\mu$, each interval $[0,s]$ with $s>0$
has positive $\mu$-mass. Thus
$A(0)-A(h)=\int_0^h\alpha(s)\dd s>0$ for every $h\in(0,1]$.
The identity $d_*(s)=A(s)+s\alpha(s)$, together with continuity
of $A$ at zero, also gives $d_*(0+)=A(0)$.
\end{proof}

For $\beta>1$, the susceptibility bound excludes the measure
$\delta_0$. This makes the temperature derivative at a fixed
minimizer strictly smaller than $\beta/2$. A finite-volume secant
estimate transfers that strict inequality to the equilibrium energy.

\begin{lemma}[Strict bound on the equilibrium energy]\label{static-lem:energy-gap}
For each fixed $\beta>1$, there are $e_0=e_0(\beta)<\beta/2$ and
$N=N(\beta)$ such that
\begin{equation}\label{static-eq:energy-gap}
 p_n'(\beta)=\frac1n\E\pi_{n,\beta}^J(H_n^J)\le e_0
 \qquad(n\ge N).
\end{equation}
\end{lemma}

\begin{proof}
Choose a minimizing measure $\mu$ at the specified $\beta$ and
hold that measure fixed during the temperature variation. Combine
Equations~\eqref{static-eq:scalar-temperature-variation} and
\eqref{static-eq:scalar-stationarity-identity}, and differentiate
the penalty, to obtain
\begin{equation}\label{static-eq:fixed-minimizer-temperature-derivative}
 \left.\frac{\dd}{\dd\gamma}\cP_\gamma(\mu)\right|_{\gamma=\beta}
 =\beta\int_0^1\bigl(d_*(s)-s\alpha(s)\bigr)\dd s
 =\beta\int_0^1 A(s)\dd s
 =\frac{\beta}{2}\left(1-\int_0^1t^2\,\mu(\dd t)\right).
\end{equation}
The last equality follows from two applications of Fubini's theorem.
If $\mu=\delta_0$, then $A(0)=1$, contradicting
$\beta A(0)\le1$ for $\beta>1$. The second moment of $\mu$ is
therefore positive. Denote the derivative in
Equation~\eqref{static-eq:fixed-minimizer-temperature-derivative}
by $L$; we have proved $L<\beta/2$.

Choose $e_0$ with $L<e_0<\beta/2$. Then fix $\eps>0$
sufficiently small that
\[
 \frac{\cP_{\beta+\eps}(\mu)-\cP_\beta(\mu)}{\eps}<e_0.
\]
For every $n$, the finite-volume pressure is differentiable and
convex in $\beta$. Bound its derivative by the forward secant,
use the scalar upper trial bound at $\beta+\eps$, and apply
Theorem~\ref{static-thm:finite-size} at $\beta$. This gives
\[
 p_n'(\beta)
 \le\frac{p_n(\beta+\eps)-p_n(\beta)}{\eps}
 \le\frac{\cP_{\beta+\eps}(\mu)-P_\beta}{\eps}
         +\frac{C_\beta}{\eps}n^{-1/24}.
\]
Since $\cP_\beta(\mu)=P_\beta$, the fixed secant slope is
strictly below $e_0$. Its strict margin absorbs the error once $n$
is sufficiently large. Finally, differentiate the finite partition
sum and then take Gaussian expectation to obtain the energy identity
in Equation~\eqref{static-eq:energy-gap}. The integrable bound
$\max_{\sigma\in\Sig}|H_n^J(\sigma)|$ justifies this passage.
The argument has varied only the temperature in the functional of
one fixed measure; it requires no differentiability of $P_\beta$
or of a choice of minimizer as a function of $\beta$.
\end{proof}
